\section{The low middle-height estimate}\label{sec:variation-estimates}

The low endpoint limits the area in one exterior quarter. As the cap gets
wider, the diagonal tent removes a quadratic amount of middle area.
We compare these two quantities, splitting only when a tent tail
reaches a boundary of the middle interval.

\begin{lemma}\label{lem:variation-low-height}
Let $U$ be a height-one downward cap with affine middle roof.
If its higher middle endpoint has height one and its lower middle
endpoint has height at most $1/2$, then, for every
$\alpha\in[\pi/4,\pi/2]$,
\[
 P_\alpha(U)<\frac45<\Mc/2.
\]
\end{lemma}

\begin{proof}
The barrier dominates the positive two-wall tent at $\pi/4$.
The resulting one-angle upper comparison is invariant under
horizontal reflection, since reflection exchanges its two diagonal
supports. We may therefore center the projection and put
\[
 I=[-2C,2C],\quad J=[-C,C],\quad A(C)=1,\quad A(-C)=a\le1/2.
\]
For $x=-2C+z$, $0\le z\le C$, concavity below the middle chord gives
\[
 A(x)\le(3a-1)/2+(1-a)z/(2C)\le1/4+z/(4C).
\]
The exterior reward is consequently at most $11C/8$.
For $C\le1/2$ this is at most $11/16$, proving the claim in
that range. Suppose $C\ge1/2$, and write
\[
 \max_U(x+y)=2C+u,\qquad\max_U(-x+y)=2C+v.
\]
The baseline endpoints, $(C,1)$, and the preceding roof bound give
\[
 \max(0,1-C)\le u\le1,\qquad0\le v\le1/4.
\]
For the last upper bound, on the left wing
$-x+A(x)\le2C+1/4-(1-1/(4C))z$.
On the middle chord it is at most $C+1/2\le2C+1/4$,
and the right wing gives an even smaller bound.

On the right wing set $z=2C-x$; then
$A(x)\le\min(1,u+z)$. Since $1-u\le C$, its integral is at most
$C-(1-u)^2/2$. On the left wing compare $v+z$ with
$1/4+z/(4C)$. Their crossing is
$z_0=(1/4-v)/(1-1/(4C))\le C$. The missing triangle gives
\[
 R\le\frac{11C}{8}-\frac{(1-u)^2}{2}-\frac{(1/4-v)^2}{2}.
\]
The actual diagonal tent is $(h-|x-d|)_+$, where
\[
 h=2C+(u+v)/2-\sqrt2,\qquad d=(u-v)/2.
\]
Its apex lies in $J$. Put $C_0=\sqrt2-1/4$. If $1/2\le C\le C_0$, its left tail
cannot reach the left endpoint, since
$h-C-d=C+v-\sqrt2\le C+1/4-\sqrt2\le0$.
Thus $C_0$ is the width at which that tail can first reach the left
endpoint; no additional width tolerance is needed.
Writing $z=C+u-\sqrt2$, exact integration therefore gives the
upper comparison
\begin{equation}\label{eq:variation-low-function}
 P_\alpha(U)\le F(C,u,v):=\frac{11C}{8}
  -\frac{(1-u)^2}{2}-\frac{(1/4-v)^2}{2}-h_+^2+\frac{z_+^2}{2}.
\end{equation}

If $z\le0$, put $x_0=1-u$, $y_0=1/4-v$ and
$H=2C+5/8-\sqrt2>0$. With $w=(x_0+y_0)/2$,
\[
 \frac{x_0^2+y_0^2}{2}+(H-w)_+^2
 \ge w^2+(H-w)_+^2\ge H^2/2.
\]
The second inequality follows by completing the square for $w\le H$,
and by $w^2\ge H^2$ for $w\ge H$.
Completing the square in $C$ now yields
\[
 F\le\frac{11\sqrt2}{16}-\frac{99}{512}
   <\frac{11}{16}\frac{10}{7}-\frac{99}{512}<4/5,
\]
using $\sqrt2<10/7$.

For $z>0$, the allowed interval in $u$ starts at $u_0=\sqrt2-C$.
Here $h>0$, and direct differentiation gives
\[
 F_u=1-C-(u+v)/2,\qquad F_{uu}=-1/2.
\]
Its maximum is either at $u_0$, already covered, or at
$u_*=2-2C-v$. The upper endpoint introduces no extra case,
since $C\ge1/2$ implies $u_*\le1$.
An interior stationary point in the clipped range satisfies
$C+v\le2-\sqrt2$. Hence
$1/4-v\ge\sqrt2-5/4>0$.
Discard the other nonnegative losses in
\eqref{eq:variation-low-function} to get
\[
 F(C,u_*,v)\le\frac{11(2-\sqrt2)}8
                     -\frac12(\sqrt2-5/4)^2
  =\frac{31}{32}-\frac{\sqrt2}{8}
  <\frac{31}{32}-\frac18\frac75<4/5,
\]
using $\sqrt2>7/5$.

Finally, for $C\ge C_0$ the baseline diagonal tent alone gives
\[
 P_\alpha(U)\le
 \begin{cases}
  11C/8-(2C-\sqrt2)^2,&C_0\le C\le\sqrt2,\\
  (11/8+2\sqrt2)C-3C^2,&C\ge\sqrt2.
 \end{cases}
\]
Both expressions decrease on their stated ranges and agree at
$C=\sqrt2$. Their largest value is therefore at the geometric
transition $C_0=\sqrt2-1/4$. At this point it is
\[
 \frac{11C_0}{8}-(2C_0-\sqrt2)^2
 =\frac{19\sqrt2}{8}-\frac{83}{32}
 <\frac{19}{8}\frac{10}{7}-\frac{83}{32}<\frac45.
\]
This bounds the remaining widths and completes the proof.
\end{proof}
